\section{Related Work}
\label{sec:related_work}

\subsection{Model Merging and Weight-Space Consolidation}
Model merging is an emerging paradigm aimed at consolidating multiple task-specific expert neural networks into a single cohesive parameter set. Early attempts such as Model Soups \cite{wortsman2022model} linearly average weights of models fine-tuned from the same initialization on identical tasks to boost generalization. For multi-task settings, Task Arithmetic \cite{ilharco2022editing} introduces task vectors that can be added or subtracted to edit model behaviors. To merge models that have undergone unconstrained fine-tuning and diverged into different loss basins, Git Re-Basin \cite{ainsworth2022git} and ZipIt! \cite{stoica2023zipit} utilize permutation matrices to align internal representation channels before averaging. More advanced approaches, including TIES-Merging \cite{yadav2023ties} and DARE \cite{yu2024dare}, explicitly address parameter interference by trimming redundant low-magnitude parameter updates, resolving sign conflicts, and dynamically scaling the remaining updates. AdaMerging \cite{yang2023adamerging} optimizes layer-wise merging coefficients on a local calibration stream to maximize multi-task accuracy. While parameter-space merging is highly elegant, it often suffers from representation collapse and task interference when experts are highly heterogeneous or when the loss landscapes are deeply non-convex.

\subsection{Dynamic Adapter Serving and Test-Time Adaptation}
Rather than consolidating expert parameters permanently in weight-space, dynamic adapter serving maintains separate, active expert paths (typically low-rank LoRA adapters \cite{hu2021lora}) on a shared, frozen backbone, and blends their representations dynamically during inference. SABLE \cite{sable2025} maps intermediate activations to ensembling weights using early-layer centroid distances, blending activation states in a single, parallel forward pass. RegCalMerge \cite{regcalmerge2025} regularizes calibration parameters during transductive serving via Class-Capacity Normalization and Elastic Spatial penalties. Sibling methods such as ChemMerge \cite{chemmerge2026} model activation flow through deep layers using non-equilibrium chemical kinetics, and TSAR \cite{tsar2026} introduces task-space anchor regularizers. At their core, these approaches rely heavily on mapping deep representations to the probability simplex using a softmax-like routing policy parameterized by temperature parameters. However, in real-world streaming deployments, the calibration split is extremely low-data (typically $16$ to $64$ samples per task). Under this transductive data scarcity, unregularized temperature optimization overfits rapidly to representation noise, leading to extreme temperature parameters, overconfident routing, and generalization collapse.

\subsection{PAC-Bayesian Generalization Theory}
PAC-Bayesian theory, pioneered by McAllester \cite{mcallester1999pac, mcallester2003simplified} and extended by Catoni \cite{catoni2007pac}, Alquier \cite{alquier2016properties}, and Guedj \cite{guedj2019primer}, provides a mathematically rigorous framework for analyzing the generalization performance of stochastic classifiers. Given a prior distribution $P$ over a hypothesis space and a calibration set of size $N$, the PAC-Bayesian bound guarantees, with high probability, that the out-of-sample risk of a learned posterior distribution $Q$ is bounded by its empirical risk plus a complexity penalty based on the Kullback-Leibler (KL) divergence $D_{\text{KL}}(Q || P)$. Traditionally, PAC-Bayesian bounds have been used to regularize stochastic neural networks by defining Gaussian priors and posteriors over the model weights. Recently, PAC-ZCA \cite{paczca2026} applied PAC-Bayesian theory to test-time adapter serving by optimizing task-specific log-temperatures under a Gaussian prior and posterior, demonstrating substantial generalization improvements over standard ERM. 

However, Gaussian modeling in PAC-ZCA is fundamentally unconstrained and ignores the fact that ensembling weights must lie on the probability simplex $\Delta^{K-1}$. By operating on unconstrained log-temperatures, Gaussian PAC-Bayes bounds are prone to numerical instability, gradient explosion, and catastrophic routing entropy collapse. In contrast, \textbf{Dirichlet-PAC} is the first framework to model ensembling weights directly on the probability simplex using a Dirichlet distribution. By deriving the exact analytical KL divergence between Dirichlet distributions, we provide the first simplex-constrained PAC-Bayesian complexity control for test-time model ensembling.
